\section{Panduan Praktis Translasi Bahasa Alami}

Translasi pernyataan bahasa alami ke rumus logika predikat sering menimbulkan kesalahan karena ambiguitas makna. Bagian ini menyajikan pedoman praktis untuk memilih pasangan kuantor dan operator secara tepat \cite{rosen2019}.

\subsection{Kamus Konstruksi Proposisi}

Tabel berikut merangkum pola translasi yang umum digunakan. Diasumsikan semesta wacana mencakup seluruh elemen relevan sesuai konteks.

\renewcommand{\arraystretch}{1.5}
\begin{longtable}{|p{4cm}|p{6cm}|p{4.5cm}|}
\caption{Kamus Konstruksi Translasi ke Logika Predikat (FOL)} \label{tab:translasi_fol} \\
\hline
\textbf{Frasa Bahasa Indonesia} & \textbf{Konversi Aturan Predikat Orde 1} & \textbf{Pola Operator Pasangan} \\ \hline
\endfirsthead

\hline
\textbf{Frasa Bahasa Indonesia} & \textbf{Konversi Aturan Predikat Orde 1} & \textbf{Pola Operator Pasangan} \\ \hline
\endhead

\multicolumn{3}{r}{\textit{berlanjut ke halaman berikutnya...}} \\
\endfoot

\hline
\endlastfoot

``Setiap entitas $P$ itu karakternya $Q$'' & $\forall x \, (P(x) \rightarrow Q(x))$ & Universal $+$ Implikasi \newline ($\forall, \rightarrow$) \\ \hline

``Beberapa / Sebagian / Ada entitas $P$ yang tergolong $Q$'' & $\exists x \, (P(x) \land Q(x))$ & Eksistensial $+$ Konjungsi \newline ($\exists, \land$) \\ \hline

``Tak satupun entitas $P$ yang tergolong $Q$'' & $\forall x \, (P(x) \rightarrow \neg Q(x))$ \newline \textit{atau} \newline $\neg \exists x \, (P(x) \land Q(x))$ & Universal $+$ negasi konsekuen \newline ($\forall, \rightarrow, \neg$) \\ \hline

``Semua entitas $P$ itu karakternya TIDAK $Q$'' & Identik / Sinonim dgn nomor urut 3 di atas & Identik \\ \hline

``Beberapa / Sebagian entitas $P$ TIDAK tergolong karakter $Q$'' & $\exists x \, (P(x) \land \neg Q(x))$ & Eksistensial $+$ $\land\neg$ \newline ($\exists, \land, \neg$) \\ \hline

``Hanya anggota $P$ murni yang boleh punya sifat $Q$'' & $\forall x \, (Q(x) \rightarrow P(x))$ \newline \textit{(*Catat: posisinya terbalik secara struktural)} & Implikasi bersilang balik \\ \hline

\end{longtable}

\subsection{Contoh Translasi Kasus Jamak}

\begin{contoh}
Diberikan domain $\mathcal{U} = \text{semua makhluk}$. Predikat dasar: \\
$A(x)$ = $x$ adalah atlet \\
$S(x)$ = $x$ berbadan sehat \\
$K(x)$ = $x$ perokok

Lakukan translasi formal:
\begin{enumerate}
    \item \textbf{Setiap atlet pasti sehat.} \\
    $\forall x \, (A(x) \rightarrow S(x))$
    \item \textbf{Tidak ada perokok yang menjadi atlet.} \\
    Frasa ``tidak ada'' mengikuti pola pada baris ketiga tabel. \\
    $\forall x \, (K(x) \rightarrow \neg A(x))$
    \item \textbf{Beberapa atlet sehat juga perokok.} \\
    Frasa ``beberapa'' menggunakan kuantor eksistensial dengan konjungsi. \\
    $\exists x \, (A(x) \land S(x) \land K(x))$
    \item \textbf{Sebagian atlet tidak sehat.} \\
    $\exists x \, (A(x) \land \neg S(x))$
    \item \textbf{Dalam organisasi ini, hanya yang tidak merokok yang sehat.} \\
    Polanya: setiap objek yang sehat harus memenuhi syarat tidak merokok. \\
    $\forall x \, (S(x) \rightarrow \neg K(x))$
\end{enumerate}
\end{contoh}
